\documentclass[12pt,a4paper,spanish]{book}
\usepackage[spanish]{babel}
\usepackage{amsmath}
\usepackage{graphicx}             
\usepackage{amsfonts}
\usepackage{amssymb}
\usepackage{graphics}
\usepackage{indentfirst}
\usepackage{enumerate}
\usepackage{fancyhdr}
\usepackage{fancybox}
\usepackage{lastpage}
\usepackage{color}
\usepackage{amsbsy}
\usepackage{amsthm}
\usepackage{listings}
\usepackage{titlesec}
\usepackage{titletoc}
\usepackage{subfig}
\renewcommand{\baselinestretch}{1}
\clubpenalty=10000
\widowpenalty=10000
\usepackage{makeidx}%glosario de términos
\usepackage[all]{xy}
\numberwithin{equation}{section}
\theoremstyle{definition}
\newtheorem{defi}{Definición}[section]

\theoremstyle{definition}
\newtheorem{propo}{Proposición}[section]

\theoremstyle{definition}
\newtheorem{coro}{Corolario}[section]

\theoremstyle{definition}
\newtheorem{teorema}{Teorema}[section]

\usepackage{srcltx}%permite la busqueda inversa, si no viene por defecto en la version de LaTeX usada.

% lo siguiente solo usarlo cuando compilo directamente en pdflatex
% si lo uso compilando en latex el {\'\i}ndice no lo hace bien
\usepackage{hyperref}

\hypersetup{colorlinks=true, linkcolor=black, urlcolor=black, citecolor=black}
\hypersetup{bookmarksopen=false, bookmarksnumbered=true} \hypersetup{pdfstartview=FitH}

\makeindex


\setlength{\parskip}{5pt plus 2pt minus 1pt}
%%%%%%%%%%%%%%%%
%%% MACROS
%%%%%%%%%%%%%%%%

\newcommand{\cplx}[1]{{{\mathcal #1}^{\scriptscriptstyle\bullet}}}

\newcommand{\dmc}[1]{{D}^-(#1)}
\newcommand{\dbc}[1]{{D}^b(#1)}

\newcommand{\fmf}[3]{{\Phi^{#1}_{{\scriptscriptstyle #2\!\rightarrow\! #3}}}}
\newcommand{\lotimes}{{\,\stackrel{\mathbf L}{\otimes}\,}}

\newcommand{\cO}{{\mathcal O}}
\newcommand{\cP}{{\mathcal P}}
\newcommand{\cE}{{\mathcal E}}
\newcommand{\cG}{{\mathcal G}}
\newcommand{\calL}{{\mathcal L}}
\newcommand{\cI}{{\mathcal I}}
\newcommand{\cF}{{\mathcal F}}
\newcommand{\cQ}{{\mathcal Q}}
\newcommand{\cM}{{\mathcal M}}

\newcommand{\f}{\mathfrak{f}}

\newcommand{\iso}{{\,\stackrel {\textstyle\sim}{\to}\,}}
\newcommand{\what}[1]{{\widehat #1}}

\newcommand{\rk}{\operatorname{rk}}
\newcommand{\ch}{\operatorname{ch}}
\newcommand{\Td}{\operatorname{Td}}
\usepackage[left=3.1cm,right=3.1cm,top=2.5cm,bottom=2.5cm]{geometry}
\begin{document}

	\begin{titlepage}

    \begin{center}

        \vspace{0.75 cm}

        {\large\textbf{Universidad de Salamanca}\\
                        Grado en Matemáticas\\}

     \vspace{2cm}
         
        \rule[1mm]{400.0pt}{1mm}

        {\Large\textbf{
            REVISIÓN DE MÉTODOS MULTIVARIANTES\\
            SUPERVISADOS Y NO SUPERVISADOS\\
            \rule[1mm]{400.0pt}{1mm}\\
        \vspace{2cm}
             Trabajo Fin de Grado}}\\

        \vspace{1cm}
        
        \begin{figure}[h!]
            \begin{center}
              \includegraphics[scale=.7]{Documentos Extra/escudo2tintas.pdf}\\
            \end{center}
		\end{figure}

		\vspace{0.75 cm}
		
        {\large\textbf{Alumno: Pedro Ángel Fraile Manzano }}

	\vspace{0.25 cm}
	{\large
	 \textbf{Tutoras: Ana Belén Nieto Librero y Nerea González García}} 
	 
	\vspace{0.5 cm}
	{\large
	 {Salamanca, Julio de 2023}} 

\end{center}

\end{titlepage}



	
	
	\pagenumbering{roman}
	\tableofcontents
	
	%\addcontentsline{toc}{chapter}{Introducción}
\chapter*{Introducción}

\section{Distribuciones Multivariantes}
\begin{defi}
Un vector aleatorio \textbf{X} es $\ldots$ %%@todo
\end{defi}
	
	
	\input{Métodos supervisados.tex}
	
	%\input{Documentos Extra/Conceptos Básicos.tex}
	
	\input{Métodos no Supervisados.tex}	
	
	
	
	
\begin{thebibliography}{99}
    \rhead[]{Bibliografía}
	
	\addcontentsline{toc}{chapter}{Bibliografía}

\bibitem{Chatfield 1989} \textbf{Chatfield,C y Collins A.J (1989)}. {\em Introduction to multivariate analysis}, Chapman and Hall.

\bibitem{Jollife 1986} \textbf{Jollife I.T.(1986)}. {\em Principal Component Analysis}, Springer-Verlag.

\bibitem{Hastie 2001} \textbf{Hastie, T.,Tibshirani, R. y Friedman J. (2001)}, \textit{The Elements of Statistical Learning, Data Mining, Inference and Prediction} Springer 

\bibitem{Cuadras 2014} \textbf{Cuadras, C.M. (2014)}, \textit{Nuevos métodos de Análisis Multivariante}, CMC Editions, Barcelona. 

\bibitem{Abdi 2010} \textbf{Abdi, H., y Williams, L. J. (2010).} \textit{Principal component analysis. Wiley Interdisciplinary Reviews: Computational Statistics}  2(4), 433–459. 
          
        
\end{thebibliography}
 


\end{document}